\centerline{\bf \Huge Сложный тургор}

\begin{flushright}
        \scriptsize{Все задачи легко гуглятся, поэтому исключительно на вашу честность!}
\end{flushright}

\problem{1}
Было 100 дверей, у каждой свой ключ (отпирающий только эту дверь). Двери пронумерованы числами $1,2, \ldots, 100$, ключи тоже, но, возможно, с ошибками: номер ключа совпадает с номером двери или отличается на 1. За одну попытку можно выбрать любой ключ, любую дверь и проверить, подходит ли этот ключ к этой двери. Можно ли гарантированно узнать, какой ключ какую дверь открывает, сделав не более

а) 99 попыток;

б) 75 попыток;

в) 74 попыток.

\begin{flushright}
        \scriptsize{1+2+3 балла}
\end{flushright}

\problem{2}
Дан правильный шестиугольник с центром $O$. Провели шесть равных окружностей с центрами в вершинах шестиугольника такие, что точка $O$ находится внутри окружностей. Угол величины $\alpha$ с вершиной $O$ высекает на этих окружностях шесть дуг. Докажите, что суммарная величина этих дуг равна $6 \alpha$.

\begin{flushright}
        \scriptsize{5 баллов}
\end{flushright}

\problem{3}
Аналитик сделал прогноз изменения курса доллара на каждый из 12 ближайших месяцев: на сколько процентов (число, большее $0 \%$ и меньшее $100 \%$ ) изменится курс за октябрь, на сколько - за ноябрь, ..., на сколько - за сентябрь. Оказалось, что про каждый месяц он верно предсказал, на сколько процентов изменится курс, но ошибся с направлением изменения (то есть, если он предсказывал, что курс увеличится на $x \%$, то курс падал на $x \%$, и наоборот). При этом через 12 месяцев курс совпал с прогнозом. В какую сторону в итоге изменился курс?

\begin{flushright}
        \scriptsize{6 баллов}
\end{flushright}


\problem{4}
Покажите, что для любой последовательности $a_0, a_1, \ldots, a_n, \ldots$, состоящей из единиц и минус единиц, найдутся такие $n$ и $k$, что

$$
\left|a_0 \cdot a_1 \cdot \ldots \cdot a_k+a_1 \cdot a_2 \cdot \ldots \cdot a_{k+1}+\ldots+a_n \cdot a_{n+1} \cdot \ldots \cdot a_{n+k}\right|=2017 .
$$

\begin{flushright}
        \scriptsize{8 баллов}
\end{flushright}

\newpage

\hspace{3mm}

\problem{5}
Кусок сыра надо разрезать на части с соблюдением таких правил: 1) вначале режем сыр на 2 куска, затем один из них режем на 2 куска, затем один из трёх кусков опять режем на 2 куска, и т.д.; 2) после каждого разрезания части могут быть разными по весу, но отношение веса любой части к весу любой другой должно быть строго больше заданного числа $R$.

a) Докажите, что при $R=0,5$ можно резать сыр так, что процесс никогда не остановится (после любого числа разрезаний можно будет отрезать ещё один кусок).

б) Докажите, что если $R>0,5$, то процесс резки когда-нибудь остановится.

в) На какое наибольшее число кусков можно разрезать сыр, если $R=0,6$ ?

\begin{flushright}
        \scriptsize{3+4+4 балла}
\end{flushright}

\problem{6}
Дан треугольник $A B C$. Пусть $I$ - центр вневписанной окружности, касающейся стороны $A B$, а $A_1$ и $B_1$ - точки касания двух других вневписанных окружностей со сторонами $B C$ и $A C$ соответственно. Пусть $M$ - середина отрезка $I C$, а отрезки $A A_1$ и $B B_1$ пересекаются в точке $N$. Докажите, что точки $N, B_1, A$ и $M$ лежат на одной окружности.

\begin{flushright}
        \scriptsize{10 баллов}
\end{flushright}

\problem{7}
Город имеет вид квадрата $n \times n$, разбитого на кварталы $1 \times 1$. Улицы идут с севера на юг и с запада на восток. Человек каждый день утром идёт из юго-западного угла в северовосточный, двигаясь только на север или восток, а вечером возвращается обратно, двигаясь только на юг или запад. Каждое утро он выбирает свой путь так, чтобы суммарная длина знакомых участков пути (тех, которые он уже проходил в том или ином направлении) была минимальна, и каждый вечер тоже. Докажите, что за $n$ дней он пройдёт все улицы целиком.

\begin{flushright}
        \scriptsize{10 баллов}
\end{flushright}